\section{Introduction}

\begin{multicols}{2}
The evolution of propulsion systems in Formula 1 represents a continuous pursuit of peak thermodynamic efficiency intertwined with the demands of highly competitive motorsport. Historically, the sport relied entirely on high-revving, naturally aspirated internal combustion engines (ICE) that prioritized peak power output over fuel efficiency. However, a sweeping paradigm shift occurred in 2014 with the introduction of the V6 turbo-hybrid era. This regulation change mandated the integration of complex Energy Recovery Systems (ERS), comprising the Motor Generator Unit-Kinetic (MGU-K) and the Motor Generator Unit-Heat (MGU-H), operating in tandem with the internal combustion engine. While this marked a monumental step toward road-relevant hybridization and record-breaking thermal efficiency, the combustion engine remained the dominant source of propulsive power, with the electrical components serving primarily as supplementary performance enhancers.

The 2026 regulatory cycle introduces a fundamental turning point in the sport's technical trajectory, decisively redefining the balance of power generation. These new regulations mandate a drastic shift toward electrical propulsion, establishing a near-equal power split between the internal combustion engine and the electrical system. The elimination of the complex MGU-H is counterbalanced by a substantial upgrade to the MGU-K, which will see its deployment capacity nearly triple to 350 kilowatts (kW). Concurrently, the energy storage system is scaled to manage an operational capacity of approximately 4 megajoules (MJ) per lap. This radical architectural restructuring dictates that the electrical powertrain transition from a supporting role to a primary driver of performance, while the ICE transitions to utilize fully sustainable drop-in fuels to further the championship's net-zero carbon objectives.

This unprecedented reliance on electrical energy introduces a myriad of profound engineering challenges. The packaging of the heavily uprated MGU-K, alongside the significantly larger and denser battery pack within the chassis, requires meticulous spatial optimization. This electrical mass significantly influences the vehicle's weight distribution, moment of inertia, and center of gravity. Furthermore, thermal management becomes an acute constraint; rapidly dissipating the immense heat flux generated by a 350 kW electrical motor and a high-cycling 4 MJ battery mandates aggressive cooling architectures that inherently conflict with aerodynamic drag reduction goals. Energy management strategies must also be entirely overhauled. The increased electrical deployment ratio necessitates highly sophisticated, algorithmic deployment and harvesting profiles to prevent severe power derating—and subsequent velocity drop-offs—at the end of long straights. Aerodynamically, the 2026 regulations mandate the implementation of active aerodynamic systems to mitigate drag and compensate for the altered power delivery characteristics, adding critical complexity to the vehicle's dynamic behavior.

The purpose of this paper is to present a comprehensive, multi-disciplinary design study of a 2026-compliant Formula 1 car. By analyzing the symbiotic relationship between the drastically altered power unit architecture and the new generative aerodynamic regulations, this study systematically explores optimal solutions for packaging, thermal management, and energy deployment. Utilizing advanced computational fluid dynamics (CFD) and transient lap time simulation (LTS), we evaluate the cascading effects of the 350 kW electrical architecture on chassis dynamics and overall vehicle concept. Ultimately, this research aims to establish a theoretical engineering framework for maximizing competitive performance under the constraints of the most electrically dominant regulatory cycle in the history of the sport.
\end{multicols}
